%═══════════════════════════════════════════════════════════════════════════════
\chapter{Modelos generativos y difusión: de la teoría a los algoritmos}
\label{ch:cap12}
%═══════════════════════════════════════════════════════════════════════════════

Los Capítulos~\ref{ch:cap10} y~\ref{ch:cap11} construyeron la maquinaria matemática completa: el proceso
OU como forward process, el teorema de Anderson como reverse SDE, y el denoising
score matching como pérdida de entrenamiento. Este capítulo ensambla esos
ingredientes en algoritmos concretos, los conecta con las arquitecturas del
Capítulo~\ref{ch:cap9}, y cierra el arco del libro con una lectura unificada de todos los
modelos de difusión como instancias del puente de Schrödinger del Capítulo~\ref{ch:cap11}.

La tensión central es aparente: DDPM, score-based SDE, DDIM y flow matching
parecen algoritmos distintos. La resolución es que todos resuelven el mismo
problema variacional —transportar $\mathcal{N}(\bm{0},\bm{I})$ a $p_0$
minimizando la KL entre leyes de procesos— con distintos niveles de
regularización entrópica $\epsilon$ y distintas elecciones de discretización.
La diferencia es de ingeniería, no de fundamento matemático.

Hay una ecuación que sintetiza el libro entero. La SDE de Langevin del
Capítulo~\ref{ch:cap5} y el reverse SDE de Anderson del Capítulo~\ref{ch:cap11} son la misma ecuación
con roles opuestos: en el SGD, el ruido isótropo destruye correlaciones para
escapar mínimos locales; en los modelos de difusión, el score reconstruye
correlaciones para generar datos. La diferencia es el signo del tiempo.

%───────────────────────────────────────────────────────────────────────────────
\section{El problema generativo y el marco unificador}
\label{sec:generativo_marco}
\dificultad{2}{1}
%───────────────────────────────────────────────────────────────────────────────

\subsection{Qué es un modelo generativo}
\label{subsec:modelo_generativo}

Dado acceso a muestras $\{\bm{x}^{(i)}\}_{i=1}^N \overset{\mathrm{iid}}{\sim}
p_0$ con $p_0$ desconocida en alta dimensión $\bm{x}\in\mathbb{R}^d$, un
\textbf{modelo generativo} produce nuevas muestras distribuidas aproximadamente
según $p_0$. Para imágenes, $d = C\times H\times W$: del orden de $10^5$--$10^6$.

Tres enfoques estructuralmente distintos:

\begin{enumerate}
  \item \textbf{Densidad explícita:} estimar $\hat{p}_0$ por máxima
        verosimilitud y muestrear. Intractable en alta dimensión salvo para
        familias paramétricas muy restringidas (flujos normalizadores).
  \item \textbf{Densidad implícita:} aprender un mapa $G_\theta:\mathcal{N}(\bm{0},\bm{I})
        \to p_0$ sin conocer $p_0$ explícitamente, usando un discriminador para
        medir la discrepancia. Ejemplo: GANs. Inestables por el equilibrio
        minimax.
  \item \textbf{Transporte:} definir un forward process conocido que lleva $p_0$
        a $\mathcal{N}(\bm{0},\bm{I})$, y aprender a invertirlo. Ejemplo:
        modelos de difusión. El objetivo de entrenamiento es convexo.
\end{enumerate}

Los modelos de difusión pertenecen al tercer enfoque. Su ventaja sobre las GANs
es estructural: el objetivo MSE sobre el ruido predicho es convexo y estable,
mientras que el equilibrio minimax de las GANs requiere un balance delicado que
fácilmente se rompe.

\subsection{El marco unificador: el puente de Schrödinger}
\label{subsec:marco_unificador}

El problema de generar muestras de $p_0$ a partir de $\mathcal{N}(\bm{0},\bm{I})$
tiene una formulación variacional precisa. Sea $\mathbb{R}$ la ley del forward
process (el proceso de referencia que lleva $p_0$ a $\mathcal{N}(\bm{0},\bm{I})$)
y $\Pi(p_T, p_0)$ el conjunto de todos los procesos de Markov con marginal
$p_T = \mathcal{N}(\bm{0},\bm{I})$ en $t=0$ y marginal $p_0$ en $t=T$.
El \textbf{puente de Schrödinger} es la solución del problema:
\[
  \mathbb{P}^* = \argmin_{\mathbb{P}\in\Pi(p_T,p_0)}
  D_{\mathrm{KL}}(\mathbb{P}\|\mathbb{R}).
\]
La Proposición~\ref{prop:anderson_schrodinger} del
Capítulo~\ref{ch:cap11} demostró que el proceso reverso de Anderson
es exactamente $\mathbb{P}^*$: el proceso más cercano al forward en KL
entre leyes de procesos, entre todos los que tienen las marginales correctas.

Lo que distingue a los distintos modelos de difusión es el nivel de
\textbf{regularización entrópica} $\epsilon$ del proceso de referencia.
Intuitivamente, $\epsilon$ mide cuánto ruido browniano contiene el proceso
de referencia $\mathbb{R}$. Cuando $\epsilon$ es grande, $\mathbb{R}$
es muy difusivo y las trayectorias son altamente estocásticas: el reverse
SDE necesita muchos pasos pequeños para invertir esa difusión (DDPM,
$\epsilon=\infty$, $\mathrm{NFE}\sim 10^3$). Cuando $\epsilon\to 0$,
el proceso de referencia degenera en un mapa determinista y las
trayectorias son rectas: el reverse ODE necesita muy pocos pasos
(OT flow matching, $\epsilon=0^+$, $\mathrm{NFE}\sim 1$--$10$).

El parámetro $\epsilon$ es la misma temperatura que aparece en la
distribución de Gibbs del Capítulo~\ref{ch:cap5} y en la escala de
CFG de la Sección~\ref{subsec:cfg}: controla el trade-off entre
diversidad estocástica y determinismo del transporte. Todos los modelos
de este capítulo son instancias del puente de Schrödinger con distintos
valores de $\epsilon$ y distintas elecciones de discretización:

\begin{table}[ht]
\centering
\begin{tabular}{lllll}
\toprule
\textbf{Modelo} & $\epsilon$ & \textbf{Proceso} & \textbf{Pérdida}
  & \textbf{NFE} \\
\midrule
DDPM      & $\infty$ & SDE discreto
  & $\|\bm{\varepsilon}-\bm{\varepsilon}_\theta\|^2$ & $10^3$ \\
Score SDE & grande   & SDE continuo
  & Score matching                                   & $10^2$ \\
DDIM      & grande   & ODE discreta
  & $\|\bm{\varepsilon}-\bm{\varepsilon}_\theta\|^2$ & $20$--$100$ \\
DPM-Solver& grande   & ODE continua
  & $\|\bm{\varepsilon}-\bm{\varepsilon}_\theta\|^2$ & $10$--$20$ \\
OT flow matching & $0^+$ & ODE continua
  & $\|v-v_\theta\|^2$                               & $1$--$10$ \\
\bottomrule
\end{tabular}
\caption{Todos los modelos de difusión como instancias del puente de
Schrödinger con distintos niveles de regularización entrópica $\epsilon$.
NFE = número de evaluaciones de la red por muestra generada.}
\label{tab:unificacion}
\end{table}

La columna de $\epsilon$ es la que el lector debe leer primero.
DDPM y score SDE no tienen trayectorias rectas porque su proceso
de referencia es muy difusivo; DDIM ya usa la ODE de flujo (determinista)
pero hereda el score network entrenado con el proceso estocástico;
OT flow matching elige directamente un proceso de referencia de difusión
nula y aprende el campo de velocidades del transporte óptimo.
El resto del capítulo desarrolla cada fila de esta tabla.

%───────────────────────────────────────────────────────────────────────────────
\section{DDPM: el modelo fundacional}
\label{sec:ddpm}
\dificultad{3}{2}
%───────────────────────────────────────────────────────────────────────────────

\subsection{El forward process discreto y su marginalización}
\label{subsec:ddpm_forward}

El forward process discreto (Ho et al., 2020) define una cadena de Markov
gaussiana con $T$ pasos y \textbf{noise schedule} $0 < \beta_1 \leq \cdots
\leq \beta_T < 1$:
\begin{equation}
  q(\bm{x}_t|\bm{x}_{t-1})
  = \mathcal{N}\!\left(\bm{x}_t;\,\sqrt{1-\beta_t}\,\bm{x}_{t-1},\,\beta_t\bm{I}\right).
  \label{eq:ddpm_forward}
\end{equation}
Definimos $\alpha_t = 1-\beta_t$ y $\bar{\alpha}_t = \prod_{s=1}^t\alpha_s$.

\begin{proposicion}{Marginalización directa del forward process}
\label{prop:ddpm_marginal}
\begin{equation}
  q(\bm{x}_t|\bm{x}_0)
  = \mathcal{N}\!\left(\bm{x}_t;\,\sqrt{\bar{\alpha}_t}\,\bm{x}_0,\,
  (1-\bar{\alpha}_t)\bm{I}\right).
  \label{eq:ddpm_marginal}
\end{equation}
\end{proposicion}

\begin{proof}
Por inducción. Para $t=1$: directo desde~\eqref{eq:ddpm_forward}.
Dado~\eqref{eq:ddpm_marginal} para $t-1$, escribir $\bm{x}_{t-1} =
\sqrt{\bar{\alpha}_{t-1}}\bm{x}_0 + \sqrt{1-\bar{\alpha}_{t-1}}\bm{\varepsilon}_{t-1}$
y $\bm{x}_t = \sqrt{\alpha_t}\bm{x}_{t-1} + \sqrt{1-\alpha_t}\bm{\varepsilon}_t$:
\[
  \bm{x}_t
  = \sqrt{\alpha_t\bar{\alpha}_{t-1}}\bm{x}_0
  + \underbrace{\sqrt{\alpha_t(1-\bar{\alpha}_{t-1})}\bm{\varepsilon}_{t-1}
  + \sqrt{1-\alpha_t}\bm{\varepsilon}_t}_{\text{gaussiano, varianza }
  \alpha_t(1-\bar{\alpha}_{t-1})+(1-\alpha_t) = 1-\bar{\alpha}_t}.
\]
Luego $\bm{x}_t = \sqrt{\bar{\alpha}_t}\bm{x}_0 + \sqrt{1-\bar{\alpha}_t}
\bm{\varepsilon}$, $\bm{\varepsilon}\sim\mathcal{N}(\bm{0},\bm{I})$. $\square$
\end{proof}

Esta reparametrización permite muestrear $\bm{x}_t$ en un único paso dado
$\bm{x}_0$ durante el entrenamiento: no es necesario iterar $t$ pasos del
forward.

Para $t\to T$ con $\bar{\alpha}_T\approx 0$: $q(\bm{x}_T|\bm{x}_0)\approx
\mathcal{N}(\bm{0},\bm{I})$, independiente de $\bm{x}_0$. Esta es la condición
de colapso al prior: el forward process destruye toda la información de $\bm{x}_0$
en $T$ pasos.

\subsection{El reverse process y el posterior analítico}
\label{subsec:ddpm_reverse_prior}

El reverse step verdadero $q(\bm{x}_{t-1}|\bm{x}_t)$ es intractable, pero el
posterior condicionado a $\bm{x}_0$ es analítico:

\begin{proposicion}{Posterior del reverse step condicionado a $\bm{x}_0$}
\label{prop:ddpm_posterior}
\begin{equation}
  q(\bm{x}_{t-1}|\bm{x}_t,\bm{x}_0)
  = \mathcal{N}\!\left(\bm{x}_{t-1};\,\tilde{\bm{\mu}}_t(\bm{x}_t,\bm{x}_0),\,
  \tilde{\beta}_t\bm{I}\right),
  \label{eq:ddpm_posterior}
\end{equation}
con:
\begin{align}
  \tilde{\bm{\mu}}_t(\bm{x}_t,\bm{x}_0)
  &= \frac{\sqrt{\bar{\alpha}_{t-1}}\beta_t}{1-\bar{\alpha}_t}\bm{x}_0
  + \frac{\sqrt{\alpha_t}(1-\bar{\alpha}_{t-1})}{1-\bar{\alpha}_t}\bm{x}_t,
  \label{eq:ddpm_posterior_mean}\\
  \tilde{\beta}_t
  &= \frac{(1-\bar{\alpha}_{t-1})\beta_t}{1-\bar{\alpha}_t}.
  \label{eq:ddpm_posterior_var}
\end{align}
\end{proposicion}

\begin{proof}
Por la regla de Bayes: $q(\bm{x}_{t-1}|\bm{x}_t,\bm{x}_0) \propto
q(\bm{x}_t|\bm{x}_{t-1})q(\bm{x}_{t-1}|\bm{x}_0)$.
Los dos factores son gaussianos en $\bm{x}_{t-1}$. Completando el cuadrado:
$\log q(\bm{x}_t|\bm{x}_{t-1}) + \log q(\bm{x}_{t-1}|\bm{x}_0)$ es
cuadrático en $\bm{x}_{t-1}$ con coeficiente del término cuadrático
$-\frac{1}{2}\left(\frac{\alpha_t}{1-\alpha_t} + \frac{1}{1-\bar{\alpha}_{t-1}}\right)
= -\frac{1}{2\tilde{\beta}_t}$
y término lineal que da la media $\tilde{\bm{\mu}}_t$. $\square$
\end{proof}

Sustituyendo la reparametrización $\bm{x}_0 = (\bm{x}_t - \sqrt{1-\bar{\alpha}_t}
\bm{\varepsilon})/\sqrt{\bar{\alpha}_t}$ en~\eqref{eq:ddpm_posterior_mean}:
\begin{equation}
  \tilde{\bm{\mu}}_t(\bm{x}_t,\bm{\varepsilon})
  = \frac{1}{\sqrt{\alpha_t}}\!\left(\bm{x}_t
  - \frac{\beta_t}{\sqrt{1-\bar{\alpha}_t}}\bm{\varepsilon}\right).
  \label{eq:ddpm_posterior_reparametrizado}
\end{equation}
El modelo reverso parametriza la media en función de $\bm{\varepsilon}$:
\begin{equation}
  \bm{\mu}_\theta(\bm{x}_t,t)
  = \frac{1}{\sqrt{\alpha_t}}\!\left(\bm{x}_t
  - \frac{\beta_t}{\sqrt{1-\bar{\alpha}_t}}\bm{\varepsilon}_\theta(\bm{x}_t,t)\right),
  \label{eq:ddpm_mu_theta}
\end{equation}
con $\bm{\varepsilon}_\theta(\bm{x}_t,t)$ una red neuronal que predice el ruido.

\subsection{Derivación completa del ELBO y la función de pérdida}
\label{subsec:ddpm_elbo}

\begin{teorema}{ELBO de DDPM y pérdida de predicción de ruido}
\label{thm:ddpm_elbo}
La cota inferior variacional de $\log p_\theta(\bm{x}_0)$ se descompone como:
\begin{equation}
  \log p_\theta(\bm{x}_0)
  \geq -L_T - \sum_{t=2}^{T} L_{t-1} - L_0,
  \label{eq:ddpm_elbo}
\end{equation}
donde:
\begin{align}
  L_T &= D_{\mathrm{KL}}\!\left(q(\bm{x}_T|\bm{x}_0)\,\|\,p(\bm{x}_T)\right),
  \label{eq:lt}\\
  L_{t-1} &= D_{\mathrm{KL}}\!\left(q(\bm{x}_{t-1}|\bm{x}_t,\bm{x}_0)\,\|\,
  p_\theta(\bm{x}_{t-1}|\bm{x}_t)\right),
  \label{eq:lt-1}\\
  L_0 &= -\mathbb{E}_{q(\bm{x}_1|\bm{x}_0)}\!\left[\log p_\theta(\bm{x}_0|\bm{x}_1)\right].
  \label{eq:l0}
\end{align}
Cada término $L_{t-1}$ entre gaussianas se reduce a:
\begin{equation}
  L_{t-1}
  = \mathbb{E}_{\bm{x}_0,\bm{\varepsilon}}\!\left[
  \frac{\beta_t^2}{2\sigma_t^2\alpha_t(1-\bar{\alpha}_t)}
  \left\|\bm{\varepsilon} - \bm{\varepsilon}_\theta(\bm{x}_t,t)\right\|^2
  \right],
  \label{eq:lt-1_simple}
\end{equation}
con $\bm{x}_t = \sqrt{\bar{\alpha}_t}\bm{x}_0 + \sqrt{1-\bar{\alpha}_t}\bm{\varepsilon}$.
La pérdida DDPM con pesos uniformes $\lambda_t = 1$ es:
\begin{equation}
  L_{\mathrm{DDPM}}(\theta)
  = \mathbb{E}_{t\sim\mathcal{U}[1,T],\,\bm{x}_0\sim p_0,\,
  \bm{\varepsilon}\sim\mathcal{N}(\bm{0},\bm{I})}\!\left[
  \left\|\bm{\varepsilon} - \bm{\varepsilon}_\theta\!\left(\sqrt{\bar{\alpha}_t}\bm{x}_0
  + \sqrt{1-\bar{\alpha}_t}\bm{\varepsilon},\,t\right)\right\|^2
  \right].
  \label{eq:ddpm_loss}
\end{equation}
\end{teorema}

\begin{proof}
\textbf{Paso 1: derivación del ELBO.}
Aplicando la desigualdad de Jensen $\log\mathbb{E}[X] \geq \mathbb{E}[\log X]$
vía la medida variacional $q(\bm{x}_{1:T}|\bm{x}_0)$:
\begin{align*}
  \log p_\theta(\bm{x}_0)
  &= \log \int p_\theta(\bm{x}_{0:T})\,d\bm{x}_{1:T}\\
  &= \log \int q(\bm{x}_{1:T}|\bm{x}_0)
    \frac{p_\theta(\bm{x}_{0:T})}{q(\bm{x}_{1:T}|\bm{x}_0)}\,d\bm{x}_{1:T}\\
  &\geq \mathbb{E}_{q(\bm{x}_{1:T}|\bm{x}_0)}\!\left[
    \log\frac{p_\theta(\bm{x}_{0:T})}{q(\bm{x}_{1:T}|\bm{x}_0)}\right].
\end{align*}
Expandiendo $p_\theta(\bm{x}_{0:T}) = p(\bm{x}_T)\prod_{t=1}^T
p_\theta(\bm{x}_{t-1}|\bm{x}_t)$ y $q(\bm{x}_{1:T}|\bm{x}_0) = \prod_{t=1}^T
q(\bm{x}_t|\bm{x}_{t-1})$, y usando la propiedad de Markov para escribir
$q(\bm{x}_t|\bm{x}_{t-1},\bm{x}_0) = q(\bm{x}_t|\bm{x}_{t-1})$:
\begin{align*}
  \text{ELBO}
  &= \mathbb{E}_q\!\left[\log\frac{p(\bm{x}_T)\prod_t p_\theta(\bm{x}_{t-1}|\bm{x}_t)}
    {\prod_t q(\bm{x}_t|\bm{x}_{t-1})}\right]\\
  &= \mathbb{E}_q\!\left[\log p(\bm{x}_T) - \log q(\bm{x}_T|\bm{x}_0)
    + \sum_{t=2}^T\log\frac{p_\theta(\bm{x}_{t-1}|\bm{x}_t)}{q(\bm{x}_{t-1}|\bm{x}_t,\bm{x}_0)}
    + \log p_\theta(\bm{x}_0|\bm{x}_1)\right]\\
  &= -L_T - \sum_{t=2}^T L_{t-1} - L_0,
\end{align*}
donde en el segundo paso se usó la regla de Bayes $q(\bm{x}_{t-1}|\bm{x}_{t-1})
= q(\bm{x}_{t-1}|\bm{x}_t,\bm{x}_0)q(\bm{x}_t|\bm{x}_0)/q(\bm{x}_{t-1}|\bm{x}_0)$
y se cancelaron los factores telescópicos $q(\bm{x}_{t-1}|\bm{x}_0)$.

\textbf{Paso 2: cálculo de $L_{t-1}$.}
Ambas distribuciones en~\eqref{eq:lt-1} son gaussianas con la misma varianza
$\sigma_t^2\bm{I}$ (donde Ho et al.\ eligen $\sigma_t^2 = \tilde{\beta}_t$ o
$\sigma_t^2 = \beta_t$, ambos válidos). La KL entre gaussianas de igual varianza
es:
\[
  D_{\mathrm{KL}}(\mathcal{N}(\bm{\mu}_1,\sigma^2\bm{I})\|\mathcal{N}(\bm{\mu}_2,\sigma^2\bm{I}))
  = \frac{1}{2\sigma^2}\|\bm{\mu}_1 - \bm{\mu}_2\|^2.
\]
Sustituyendo $\bm{\mu}_1 = \tilde{\bm{\mu}}_t(\bm{x}_t,\bm{x}_0)$ de~\eqref{eq:ddpm_posterior_mean}
y $\bm{\mu}_2 = \bm{\mu}_\theta(\bm{x}_t,t)$ de~\eqref{eq:ddpm_mu_theta}:
\begin{align*}
  \|\tilde{\bm{\mu}}_t - \bm{\mu}_\theta\|^2
  &= \left\|\frac{1}{\sqrt{\alpha_t}}\!\left(\bm{x}_t
    - \frac{\beta_t}{\sqrt{1-\bar{\alpha}_t}}\bm{\varepsilon}\right)
    - \frac{1}{\sqrt{\alpha_t}}\!\left(\bm{x}_t
    - \frac{\beta_t}{\sqrt{1-\bar{\alpha}_t}}\bm{\varepsilon}_\theta\right)\right\|^2\\
  &= \frac{\beta_t^2}{\alpha_t(1-\bar{\alpha}_t)}
    \|\bm{\varepsilon} - \bm{\varepsilon}_\theta\|^2.
\end{align*}
Dividiendo por $2\sigma_t^2$ se obtiene~\eqref{eq:lt-1_simple}.

\textbf{Paso 3: simplificación de la pérdida.}
El prefactor $\frac{\beta_t^2}{2\sigma_t^2\alpha_t(1-\bar{\alpha}_t)}$ varía con
$t$. Ho et al.\ observan empíricamente que omitir los pesos (tomar $\lambda_t = 1$)
mejora la calidad de las muestras: los términos con $t$ pequeño (alta resolución)
son numéricamente pequeños con los pesos correctos pero contribuyen
desproporcionadamente a la calidad perceptual. La pérdida simplificada es
entonces la esperanza sobre $t$ uniforme de $\|\bm{\varepsilon} -
\bm{\varepsilon}_\theta(\bm{x}_t,t)\|^2$, que es exactamente~\eqref{eq:ddpm_loss}.
$\square$
\end{proof}

\begin{observacion}{Conexión con el denoising score matching del Capítulo~\ref{ch:cap11}}
La pérdida~\eqref{eq:ddpm_loss} es idéntica a la pérdida de denoising score matching
del Capítulo~\ref{ch:cap11} derivada por Vincent (2011).
Esto no es coincidencia: el ELBO y el score matching son dos rutas hacia la
misma función de pérdida. El ELBO usa la estructura de la cadena de Markov y la
KL entre gaussianas; el score matching usa la dualidad $L^2$ del generador.
Ambos convergen a predecir el ruido $\bm{\varepsilon}$.
\end{observacion}

\begin{fisicaconexion}{El ELBO en física estadística y la energía libre variacional}
El ELBO (Evidence Lower Bound) del Teorema~\ref{thm:ddpm_elbo} tiene una
interpretación directa en física estadística. Para una distribución de Boltzmann
$p^*(\bm{x}) \propto e^{-E(\bm{x})/(k_BT)}$ con energía $E$ intractable, se
aproxima con una distribución variacional $q_\phi(\bm{x})$ (gaussiana u otra
familia tratable). La desigualdad de Jensen implica:
\[
  \log Z = \log\int e^{-E/(k_BT)}\,d\bm{x}
  \geq \mathbb{E}_{q_\phi}\!\left[-\frac{E}{k_BT}\right] - D_\mathrm{KL}(q_\phi\|p^*),
\]
que es exactamente el ELBO: maximizarlo minimiza la \textbf{energía libre
variacional} $F_\phi = \langle E\rangle_{q_\phi} - T S(q_\phi)$ (energía media
menos entropía de la aproximación). En el DDPM, el rol de $E$ lo juega
$-\log p_\theta(\bm{x}_{0:T})$ y el rol de $q$ lo juega la distribución forward
$q(\bm{x}_{1:T}|\bm{x}_0)$: el ELBO del DDPM es la energía libre variacional
del sistema de difusión inversa.

Los términos $L_T$, $L_{t-1}$ y $L_0$ en~\eqref{eq:ddpm_elbo} corresponden
respectivamente a: (1) la energía libre del estado terminal (cuánto difiere el
prior del estado final del forward), (2) la entropía de reconstrucción a cada
paso (la KL entre el posterior verdadero y el aprendido), y (3) la energía libre
del primer paso (calidad de reconstrucción desde $\bm{x}_1$). La simplificación
que ignora los pesos temporales es análoga a la termodinámica de temperatura
uniforme: tratar todos los pasos de la cadena como contribuyentes igualmente a
la energía libre total.

En los \textbf{autoencoders variacionales} (VAEs, Capítulo~\ref{ch:cap7}), el
mismo ELBO aparece como $\mathcal{L}(\theta,\phi) = \mathbb{E}_{q_\phi(\bm{z}|\bm{x})}
[\log p_\theta(\bm{x}|\bm{z})] - D_\mathrm{KL}(q_\phi(\bm{z}|\bm{x})\|p(\bm{z}))$:
el primer término es la verosimilitud de reconstrucción y el segundo es el
regularizador que acerca la distribución latente al prior. El VAE es un DDPM
con $T=1$ paso y sin el factor de difusión progresiva.
\end{fisicaconexion}

\subsection{El algoritmo DDPM}
\label{subsec:ddpm_algoritmo}

\begin{algorithm}[H]
\caption{DDPM: entrenamiento y muestreo}
\label{alg:ddpm}
\SetKwBlock{Train}{Entrenamiento:}{}
\SetKwBlock{Sample}{Muestreo:}{}
\Train{
  \Repeat{convergencia}{
    $\bm{x}_0 \sim p_0$ (muestra del dataset)\;
    $t \sim \mathcal{U}\{1,\ldots,T\}$\;
    $\bm{\varepsilon} \sim \mathcal{N}(\bm{0},\bm{I})$\;
    $\bm{x}_t \leftarrow \sqrt{\bar{\alpha}_t}\bm{x}_0 + \sqrt{1-\bar{\alpha}_t}\bm{\varepsilon}$\;
    Actualizar $\theta$ con gradiente de $\|\bm{\varepsilon} - \bm{\varepsilon}_\theta(\bm{x}_t,t)\|^2$\;
  }
}
\Sample{
  $\bm{x}_T \sim \mathcal{N}(\bm{0},\bm{I})$\;
  \For{$t = T, T-1, \ldots, 1$}{
    $\bm{z} \sim \mathcal{N}(\bm{0},\bm{I})$ si $t > 1$, sino $\bm{z} = \bm{0}$\;
    $\bm{x}_{t-1} \leftarrow \frac{1}{\sqrt{\alpha_t}}\!\left(\bm{x}_t -
    \frac{\beta_t}{\sqrt{1-\bar{\alpha}_t}}\bm{\varepsilon}_\theta(\bm{x}_t,t)\right)
    + \sigma_t\bm{z}$\;
  }
  \Return $\bm{x}_0$\;
}
\end{algorithm}

\subsection{El noise schedule y su conexión con el proceso OU}
\label{subsec:ddpm_schedule}

El schedule lineal de Ho et al.: $\beta_t = \beta_{\min} + (t-1)/T\cdot
(\beta_{\max}-\beta_{\min})$ con $\beta_{\min}=10^{-4}$, $\beta_{\max}=0.02$.
El schedule coseno (Nichol \& Dhariwal, 2021):
$\bar{\alpha}_t = \cos^2\!\left(\frac{t/T + s}{1+s}\cdot\frac{\pi}{2}\right)$
con $s = 0.008$, que evita la caída brusca de señal-a-ruido al final del
forward.

La conexión con el proceso OU continuo del Capítulo~\ref{ch:cap10} es directa: para
$\beta_t = 2\alpha\Delta t$ y $T\to\infty$, el forward discreto~\eqref{eq:ddpm_forward}
converge al proceso OU $d\bm{X}_t = -\alpha\bm{X}_t\,dt + \sqrt{2\alpha}\,
d\bm{W}_t$. En ese límite: $\bar{\alpha}_t = (1-\beta)^t \approx e^{-2\alpha t}$
(la atenuación exponencial del Capítulo~\ref{ch:cap10}) y $1-\bar{\alpha}_t \approx
1-e^{-2\alpha t} = \sigma_t^2$ (la varianza acumulada del proceso OU).
El schedule coseno es la discretización que mantiene la señal-a-ruido
$\sqrt{\bar{\alpha}_t/(1-\bar{\alpha}_t)}$ decreciendo a tasa aproximadamente
constante, lo que corresponde a una parametrización uniforme del tiempo en el
espacio logarítmico del SNR.

%───────────────────────────────────────────────────────────────────────────────
\section{Score-based SDE: el marco continuo}
\label{sec:score_sde}
\dificultad{3}{3}
%───────────────────────────────────────────────────────────────────────────────

\subsection{Formulación unificada de Song et al.\ (2021)}
\label{subsec:score_sde_formulacion}

El forward process continuo $t\in[0,T]$ con drift $\bm{f}(\bm{x},t)$ y difusión
escalar $g(t)$:
\begin{equation}
  d\bm{X}_t = \bm{f}(\bm{X}_t,t)\,dt + g(t)\,d\bm{W}_t.
  \label{eq:score_sde_forward}
\end{equation}
Dos familias canónicas que generalizan DDPM:

\textbf{VP-SDE} (variance preserving, equivalente a DDPM en el límite continuo):
$\bm{f}(\bm{x},t) = -\frac{1}{2}\beta(t)\bm{x}$, $g(t) = \sqrt{\beta(t)}$.
Las marginales son gaussianas: $p_t(\bm{x}|\bm{x}_0) = \mathcal{N}(\bm{x};
e^{-\frac{1}{2}\int_0^t\beta(s)ds}\bm{x}_0, (1-e^{-\int_0^t\beta(s)ds})\bm{I})$.
La varianza total $\mathrm{Var}[\bm{X}_t] \to 1$ para $t\to T$ grande.

\textbf{VE-SDE} (variance exploding): $\bm{f}(\bm{x},t) = \bm{0}$, $g(t) =
\sqrt{d\sigma^2(t)/dt}$ con $\sigma(t)$ creciente. La varianza explota:
$\mathrm{Var}[\bm{X}_t] = \sigma(t)^2 \to\infty$. Útil cuando se quiere
controlar explícitamente el nivel de ruido $\sigma$.

El reverse SDE del Corolario~\ref{cor:anderson_isotropo}:
\begin{equation}
  d\tilde{\bm{X}}_t = \left[-\bm{f}(\tilde{\bm{X}}_t,T{-}t)
  + g(T{-}t)^2\,\nabla_{\bm{x}}\log p_{T-t}(\tilde{\bm{X}}_t)\right]dt
  + g(T{-}t)\,d\tilde{\bm{W}}_t.
  \label{eq:score_sde_reverse}
\end{equation}
La red $\bm{s}_\theta(\bm{x},t) \approx \nabla\log p_t(\bm{x})$ se entrena
con la pérdida de denoising score matching~\eqref{eq:ddpm_loss} del Capítulo~\ref{ch:cap11}
(con la identificación $\bm{\varepsilon}_\theta = -\sigma_t\bm{s}_\theta$).

\subsection{Samplers numéricos para el reverse SDE}
\label{subsec:score_sde_samplers}

La discretización de Euler-Maruyama del reverse SDE con paso $\Delta t = T/N$:
\begin{equation}
  \bm{x}_{t-\Delta t}
  = \bm{x}_t + \left[\bm{f}(\bm{x}_t,t) - g(t)^2\bm{s}_\theta(\bm{x}_t,t)\right]\Delta t
  + g(t)\sqrt{\Delta t}\,\bm{z}, \quad \bm{z}\sim\mathcal{N}(\bm{0},\bm{I}).
  \label{eq:euler_reverse}
\end{equation}
El predictor-corrector de Song et al.: alterna~\eqref{eq:euler_reverse} (predictor)
con pasos de Langevin MCMC (corrector):
\begin{equation}
  \bm{x} \leftarrow \bm{x} + \delta\,\bm{s}_\theta(\bm{x},t) + \sqrt{2\delta}\,\bm{z},
  \label{eq:langevin_corrector}
\end{equation}
durante $M$ iteraciones por nivel de ruido. El corrector reduce el error de
discretización acumulado sin aumentar la evaluación total de la red.

\subsection{DDIM: derivación completa como ODE de flujo discretizada}
\label{subsec:ddim}

DDIM (Song et al., 2020) es el sampler que elimina el ruido del reverse step
reemplazando la SDE estocástica por la ODE de flujo de probabilidad del
Capítulo~\ref{ch:cap11}. La derivación parte de la forma general del step DDPM y muestra
que el caso determinista es un caso límite consistente.

\begin{teorema}{Paso de actualización DDIM}
\label{thm:ddim}
Dado $\bm{x}_t$ y el predictor de ruido $\bm{\varepsilon}_\theta(\bm{x}_t,t)$,
el paso DDIM hacia $\bm{x}_{t'}$ con $t' < t$ es:
\begin{equation}
  \bm{x}_{t'}
  = \sqrt{\bar{\alpha}_{t'}}\underbrace{\frac{\bm{x}_t - \sqrt{1-\bar{\alpha}_t}
  \bm{\varepsilon}_\theta(\bm{x}_t,t)}{\sqrt{\bar{\alpha}_t}}}_{\hat{\bm{x}}_0(\bm{x}_t,t)}
  + \sqrt{1-\bar{\alpha}_{t'} - \eta^2\tilde{\beta}_{t'}}\,
  \bm{\varepsilon}_\theta(\bm{x}_t,t)
  + \eta\sqrt{\tilde{\beta}_{t'}}\,\bm{z},
  \label{eq:ddim_step}
\end{equation}
donde $\bm{z}\sim\mathcal{N}(\bm{0},\bm{I})$ y $\eta\in[0,1]$ controla la
estocasticidad. Para $\eta=0$: paso completamente determinista (ODE de flujo).
Para $\eta=1$: paso DDPM estocástico.
\end{teorema}

\begin{proof}
\textbf{Paso 1: familia general de distribuciones de transición consistentes.}
La condición de consistencia marginal requiere que $q(\bm{x}_{t'}|\bm{x}_0)$
sea la gaussiana correcta~\eqref{eq:ddpm_marginal}. Consideramos la familia de
distribuciones de transición generalizadas:
\[
  q_\sigma(\bm{x}_{t'}|\bm{x}_t,\bm{x}_0)
  = \mathcal{N}\!\left(\bm{x}_{t'};\,\tilde{\bm{\mu}}_{t'}(\bm{x}_t,\bm{x}_0),\,
  \sigma^2\bm{I}\right),
\]
con media a determinar.

\textbf{Paso 2: la media correcta desde las marginales.}
Queremos $\int q_\sigma(\bm{x}_{t'}|\bm{x}_t,\bm{x}_0)q(\bm{x}_t|\bm{x}_0)
\,d\bm{x}_t = q(\bm{x}_{t'}|\bm{x}_0)$.
Usando la reparametrización $\bm{x}_t = \sqrt{\bar{\alpha}_t}\bm{x}_0 +
\sqrt{1-\bar{\alpha}_t}\bm{\varepsilon}$ y $\bm{x}_{t'} = \sqrt{\bar{\alpha}_{t'}}
\bm{x}_0 + \sqrt{1-\bar{\alpha}_{t'}}\bm{\varepsilon}'$, la media que satisface
la consistencia marginal es:
\begin{align*}
  \tilde{\bm{\mu}}_{t'}(\bm{x}_t,\bm{x}_0)
  &= \sqrt{\bar{\alpha}_{t'}}\bm{x}_0
  + \sqrt{1-\bar{\alpha}_{t'}-\sigma^2}\cdot\frac{\bm{x}_t
  - \sqrt{\bar{\alpha}_t}\bm{x}_0}{\sqrt{1-\bar{\alpha}_t}}.
\end{align*}
Esta es la media que produce exactamente $q(\bm{x}_{t'}|\bm{x}_0) =
\mathcal{N}(\sqrt{\bar{\alpha}_{t'}}\bm{x}_0,(1-\bar{\alpha}_{t'})\bm{I})$
al marginalizar sobre $q(\bm{x}_t|\bm{x}_0)$, independientemente de $\sigma$.

\textbf{Paso 3: sustitución del predictor.}
Se estima $\bm{x}_0$ a partir de $\bm{x}_t$ usando el predictor de ruido:
$\hat{\bm{x}}_0 = (\bm{x}_t - \sqrt{1-\bar{\alpha}_t}\bm{\varepsilon}_\theta)/
\sqrt{\bar{\alpha}_t}$.
Sustituyendo en $\tilde{\bm{\mu}}_{t'}$:
\[
  \tilde{\bm{\mu}}_{t'}(\bm{x}_t)
  = \sqrt{\bar{\alpha}_{t'}}\hat{\bm{x}}_0
  + \sqrt{1-\bar{\alpha}_{t'}-\sigma^2}
  \cdot\bm{\varepsilon}_\theta(\bm{x}_t,t).
\]

\textbf{Paso 4: el paso completo con parámetro $\eta$.}
Tomando $\sigma^2 = \eta^2\tilde{\beta}_{t'}$ y añadiendo el término de ruido
$\sigma\bm{z}$ se obtiene exactamente~\eqref{eq:ddim_step}.

\textbf{Paso 5: el caso $\eta=0$ como ODE de flujo discretizada.}
Para $\eta=0$, el paso es determinista:
\[
  \bm{x}_{t'} = \sqrt{\bar{\alpha}_{t'}}\hat{\bm{x}}_0
  + \sqrt{1-\bar{\alpha}_{t'}}\,\bm{\varepsilon}_\theta(\bm{x}_t,t).
\]
Reescribiendo en términos de $\bm{x}_t$:
\begin{align*}
  \bm{x}_{t'}
  &= \sqrt{\bar{\alpha}_{t'}}\cdot\frac{\bm{x}_t - \sqrt{1-\bar{\alpha}_t}
    \bm{\varepsilon}_\theta}{\sqrt{\bar{\alpha}_t}}
    + \sqrt{1-\bar{\alpha}_{t'}}\,\bm{\varepsilon}_\theta \\
  &= \sqrt{\frac{\bar{\alpha}_{t'}}{\bar{\alpha}_t}}\bm{x}_t
    + \left(\sqrt{1-\bar{\alpha}_{t'}}
    - \sqrt{\frac{\bar{\alpha}_{t'}(1-\bar{\alpha}_t)}{\bar{\alpha}_t}}\right)
    \bm{\varepsilon}_\theta.
\end{align*}
Para pasos pequeños $\Delta = t - t' \to 0$:
$\sqrt{\bar{\alpha}_{t'}/\bar{\alpha}_t} \approx 1 - \frac{1}{2}\Delta\beta(t)$ y
$\sqrt{1-\bar{\alpha}_{t'}} \approx \sqrt{1-\bar{\alpha}_t} +
\frac{\beta(t)\Delta}{2\sqrt{1-\bar{\alpha}_t}}$.
Luego:
\[
  \frac{\bm{x}_{t'} - \bm{x}_t}{\Delta t}
  \approx -\frac{\beta(t)}{2}\bm{x}_t
  + \frac{\beta(t)}{2\sqrt{1-\bar{\alpha}_t}}\bm{\varepsilon}_\theta(\bm{x}_t,t)
  = \bm{f}(\bm{x}_t,t) - \frac{g(t)^2}{2}\bm{s}_\theta(\bm{x}_t,t),
\]
que es exactamente la ODE de flujo de probabilidad~\eqref{eq:prob_flow_ode} del
Capítulo~\ref{ch:cap11} discretizada por Euler. DDIM con $\eta=0$ es la discretización de
Euler de la ODE de flujo de probabilidad. $\square$
\end{proof}

La demostración explica por qué DDIM puede generar muestras de calidad comparable
a DDPM con $20$--$50$ pasos en lugar de $1000$: la ODE de flujo tiene soluciones
más suaves que la SDE estocástica, y los solvers de EDOs de orden superior
convergen mucho más rápido.

\subsection{DPM-Solver: integración exacta del término lineal}
\label{subsec:dpm_solver}

DPM-Solver (Lu et al., 2022) explota la estructura de la ODE de flujo VP para
integrar exactamente el término lineal, dejando solo el término no lineal para
aproximación numérica.

En términos del tiempo logarítmico del SNR $\lambda_t = \frac{1}{2}\log
(\bar{\alpha}_t/(1-\bar{\alpha}_t))$, la ODE de flujo VP se escribe:
\begin{equation}
  \frac{d\bm{x}}{d\lambda}
  = \bm{x} - e^{-\lambda}(1+e^{2\lambda})\bm{s}_\theta(\bm{x},t(\lambda)),
  \label{eq:dpm_ode}
\end{equation}
donde el lado derecho tiene un término lineal exactamente integrable.
La fórmula de variación de parámetros da:
\[
  \bm{x}_{\lambda'} = e^{\lambda'-\lambda}\bm{x}_\lambda
  - \int_\lambda^{\lambda'} e^{\lambda'-r}(1+e^{2r})\bm{s}_\theta(\bm{x}(r),r)\,dr.
\]
Aproximando $\bm{s}_\theta(\bm{x}(r),r) \approx \bm{s}_\theta(\bm{x}_\lambda,\lambda)$
(DPM-Solver-1, orden 1) o con correcciones de orden superior por expansión de Taylor
en $\lambda$, se obtienen solvers de orden 1, 2 y 3. El solver de orden 3 reduce
los pasos necesarios a $10$--$20$ con error despreciable.

%───────────────────────────────────────────────────────────────────────────────
\section{Flow matching: el límite de transporte óptimo}
\label{sec:flow_matching}
\dificultad{3}{2}
%───────────────────────────────────────────────────────────────────────────────

\subsection{Interpolación y campo vectorial condicional}
\label{subsec:fm_campo}

El flow matching (Lipman et al., 2022; Liu et al., 2022) aprende el campo vectorial
$v_\theta(\bm{x},t)$ de la ODE:
\begin{equation}
  \frac{d\bm{x}}{dt} = v_\theta(\bm{x},t),
  \quad \bm{x}(0)\sim p_T = \mathcal{N}(\bm{0},\bm{I}),
  \quad \bm{x}(1)\sim p_0,
  \label{eq:fm_ode}
\end{equation}
que transporta $p_T$ a $p_0$ en tiempo $t\in[0,1]$. La pérdida de flow matching
condicional (CFM):
\begin{equation}
  L_{\mathrm{CFM}}(\theta)
  = \mathbb{E}_{t,\bm{x}_0,\bm{x}_T}\!\left[
  \left\|v_\theta\!\left((1-t)\bm{x}_T + t\bm{x}_0,\,t\right)
  - (\bm{x}_0 - \bm{x}_T)\right\|^2\right],
  \label{eq:cfm_loss}
\end{equation}
donde $\bm{x}(t) = (1-t)\bm{x}_T + t\bm{x}_0$ es la interpolación lineal y
$\bm{x}_0 - \bm{x}_T$ es la velocidad objetivo.

\begin{proposicion}{Equivalencia entre CFM y FM marginal}
\label{prop:cfm_fm}
Los objetivos de flow matching marginal $L_{\mathrm{FM}}$ (sobre el campo
$v_t^{\mathrm{marg}}(\bm{x}) = \mathbb{E}[\bm{x}_0-\bm{x}_T|\bm{x}(t)=\bm{x}]$)
y condicional $L_{\mathrm{CFM}}$ tienen el mismo gradiente respecto a $\theta$:
$\nabla_\theta L_{\mathrm{FM}} = \nabla_\theta L_{\mathrm{CFM}}$.
\end{proposicion}

\begin{proof}
$L_{\mathrm{FM}} = \mathbb{E}_{t,\bm{x}\sim p_t}\|v_\theta(\bm{x},t) -
v_t^{\mathrm{marg}}(\bm{x})\|^2$ y $L_{\mathrm{CFM}} = \mathbb{E}_{t,\bm{x}_0,\bm{x}_T}
\|v_\theta(\bm{x}(t),t) - (\bm{x}_0-\bm{x}_T)\|^2$.
Expandiendo $L_{\mathrm{CFM}} = \mathbb{E}\|v_\theta\|^2 - 2\mathbb{E}[v_\theta\cdot
(\bm{x}_0-\bm{x}_T)] + C$.
El término cruzado, por la ley de la esperanza total condicionando en
$\bm{x}(t) = \bm{x}$:
$\mathbb{E}[v_\theta(\bm{x}(t),t)\cdot(\bm{x}_0-\bm{x}_T)]
= \mathbb{E}_{t,\bm{x}}\left[v_\theta(\bm{x},t)\cdot\mathbb{E}[\bm{x}_0-\bm{x}_T|
\bm{x}(t)=\bm{x}]\right]
= \mathbb{E}_{t,\bm{x}}[v_\theta\cdot v_t^{\mathrm{marg}}]$.
Luego $L_{\mathrm{CFM}} = L_{\mathrm{FM}} + C'$ y los gradientes coinciden.
$\square$
\end{proof}

\subsection{Rectified flow e iteraciones de reflow}
\label{subsec:rectified_flow}

El \textbf{rectified flow} (Liu et al., 2022) itera el flow matching usando el
pareo inducido por el flujo aprendido en la iteración anterior:

\begin{algorithm}[H]
\caption{Rectified flow con $k$ iteraciones de reflow}
\label{alg:rectified_flow}
$\bm{x}_T^{(0)} \sim p_T$, $\bm{x}_0^{(0)} \sim p_0$ (pareo independiente)\;
\For{$k = 0, 1, \ldots, K-1$}{
  Entrenar $v_\theta^{(k)}$ con la pérdida~\eqref{eq:cfm_loss} usando el pareo
  $(\bm{x}_T^{(k)}, \bm{x}_0^{(k)})$\;
  Para cada $\bm{x}_T^{(k)}$: integrar $d\bm{x}/dt = v_\theta^{(k)}(\bm{x},t)$
  para obtener $\bm{x}_0^{(k+1)}$\;
  Actualizar pareo: $(\bm{x}_T^{(k+1)}, \bm{x}_0^{(k+1)})
  \leftarrow (\bm{x}_T^{(k)}, \bm{x}_0^{(k+1)})$\;
}
\end{algorithm}

Con cada iteración, las trayectorias se enderezan: el pareo se aproxima al
mapa de transporte óptimo $T^*$ que hace las trayectorias rectas (sección sin
cruce). La convergencia es geométrica: el costo de cruce de trayectorias
decrece exponencialmente en $k$.

\subsection{OT flow matching: pareo óptimo directo}
\label{subsec:ot_fm}

El \textbf{OT flow matching} (Tong et al., 2023) usa el plan de transporte
óptimo de Sinkhorn (Algoritmo~\ref{alg:sinkhorn}, Capítulo~\ref{ch:cap11}) como pareo
desde la primera iteración, eliminando la necesidad de reflow.

Con pareo óptimo $(\bm{x}_0, T^*(\bm{x}_0))$, las trayectorias no se cruzan y
el campo marginal satisface $v_t^{\mathrm{marg}}(\bm{x}) = T^*(\bm{x}_0(t)) -
\bm{x}_0(t)$, que es constante a lo largo de cada trayectoria. El modelo aprende
un campo más simple, converge con menos datos, y el muestreo necesita menos
pasos. Esta es la arquitectura de Stable Diffusion 3, Flux y Stable Audio 2:
todos usan flow matching con pareo por transporte óptimo.

%───────────────────────────────────────────────────────────────────────────────
\section{Condicionamiento y Classifier-Free Guidance}
\label{sec:cfg}
\dificultad{2}{2}
%───────────────────────────────────────────────────────────────────────────────

\subsection{El score condicionado: derivación vía regla de Bayes}
\label{subsec:score_condicionado}

Sea $y$ una condición (clase, texto, imagen de referencia). El objetivo es
generar $\bm{x}\sim p_0(\bm{x}|y)$. El score condicionado satisface una
identidad exacta:

\begin{proposicion}{Descomposición del score condicionado}
\label{prop:score_condicionado}
\begin{equation}
  \nabla_{\bm{x}}\log p_t(\bm{x}|y)
  = \nabla_{\bm{x}}\log p_t(\bm{x})
  + \nabla_{\bm{x}}\log p_t(y|\bm{x}).
  \label{eq:score_condicionado}
\end{equation}
\end{proposicion}

\begin{proof}
Por la regla de Bayes: $p_t(\bm{x}|y) = p_t(\bm{x})p_t(y|\bm{x})/p_t(y)$.
Tomando logaritmo y gradiente respecto a $\bm{x}$:
$\nabla\log p_t(\bm{x}|y) = \nabla\log p_t(\bm{x}) + \nabla\log p_t(y|\bm{x})
- \nabla\log p_t(y)$.
El término $\nabla\log p_t(y)$ no depende de $\bm{x}$: es cero. $\square$
\end{proof}

El reverse SDE condicionado es~\eqref{eq:score_sde_reverse} con el score
condicionado~\eqref{eq:score_condicionado}:
\[
  d\tilde{\bm{X}}_t = \left[-\bm{f} + g^2(\nabla\log p_{T-t}
  + \nabla\log p_{T-t}(y|\tilde{\bm{X}}_t))\right]dt + g\,d\tilde{\bm{W}}_t.
\]
El segundo término es el gradiente del clasificador $p_t(y|\bm{x})$ respecto a
la entrada: \textbf{classifier guidance} (Dhariwal \& Nichol, 2021).

\subsection{Classifier guidance con escala de temperatura}
\label{subsec:classifier_guidance}

El classifier guidance amplifica la condición por un factor $s > 0$:
\begin{equation}
  \nabla\log p_t^{(s)}(\bm{x}|y)
  = \nabla\log p_t(\bm{x}) + s\,\nabla\log p_t(y|\bm{x}).
  \label{eq:classifier_guidance}
\end{equation}

\begin{proposicion}{Interpretación de temperatura del classifier guidance}
\label{prop:guidance_temperatura}
La distribución $p_t^{(s)}(\bm{x}|y)$ definida por el score~\eqref{eq:classifier_guidance}
es proporcional a:
\begin{equation}
  p_t^{(s)}(\bm{x}|y) \propto p_t(\bm{x})\,p_t(y|\bm{x})^s.
  \label{eq:guidance_distribución}
\end{equation}
\end{proposicion}

\begin{proof}
Integrando el score~\eqref{eq:classifier_guidance}:
$\log p_t^{(s)}(\bm{x}|y) = \log p_t(\bm{x}) + s\log p_t(y|\bm{x}) + C$,
luego $p_t^{(s)} \propto p_t(\bm{x})p_t(y|\bm{x})^s$. $\square$
\end{proof}

La distribución~\eqref{eq:guidance_distribución} es exactamente la distribución
de Gibbs del Capítulo~\ref{ch:cap5} con temperatura $T_s = 1/s$ para el ``potencial''
$-\log p_t(y|\bm{x})$:
\[
  p_t^{(s)}(\bm{x}|y) \propto p_t(\bm{x})\,e^{s\log p_t(y|\bm{x})}
  = p_t(\bm{x})\,e^{-U(\bm{x})/T_s}, \quad U(\bm{x}) = -\log p_t(y|\bm{x}).
\]
Para $s=1$: distribución Bayesiana exacta $p(\bm{x}|y)$.
Para $s>1$: temperatura baja, distribución concentrada alrededor del máximo de
$p(y|\bm{x})$ (alta ``fidelidad'' a la condición, baja diversidad).
Para $s<1$: temperatura alta, distribución difusa (baja fidelidad, alta diversidad).
El parámetro $s$ es el control de temperatura del Capítulo~\ref{ch:cap5} en el contexto del
muestreo condicionado.

\subsection{Classifier-Free Guidance (CFG)}
\label{subsec:cfg}

El classifier guidance requiere un clasificador externo $p_t(y|\bm{x})$
entrenado con datos ruidosos $\bm{x}_t$. El \textbf{Classifier-Free Guidance}
(Ho \& Salimans, 2022) elimina el clasificador externo reescribiendo el score
condicionado en términos del score no condicionado y el score condicionado de
la red de difusión.

\begin{proposicion}{CFG como diferencia de scores}
\label{prop:cfg}
El score condicionado con guía de escala $s$ es:
\begin{equation}
  \nabla\log p_t^{(s)}(\bm{x}|y)
  = (1-s)\nabla\log p_t(\bm{x}) + s\,\nabla\log p_t(\bm{x}|y)
  + (s-1)[\nabla\log p_t(\bm{x}|y) - \nabla\log p_t(\bm{x})].
  \label{eq:cfg_score}
\end{equation}
Equivalentemente, en términos del predictor de ruido:
\begin{equation}
  \tilde{\bm{\varepsilon}}_\theta(\bm{x}_t,y,t)
  = \bm{\varepsilon}_\theta(\bm{x}_t,t)
  + s\left[\bm{\varepsilon}_\theta(\bm{x}_t,y,t)
    - \bm{\varepsilon}_\theta(\bm{x}_t,t)\right],
  \label{eq:cfg_eps}
\end{equation}
donde $\bm{\varepsilon}_\theta(\bm{x}_t,t)$ es la predicción no condicionada
(obtenida entrenando con $y = \emptyset$ con probabilidad $p_{\mathrm{drop}}$)
y $\bm{\varepsilon}_\theta(\bm{x}_t,y,t)$ es la predicción condicionada.
\end{proposicion}

\begin{proof}
De~\eqref{eq:score_condicionado}: $\nabla\log p_t(y|\bm{x}) = \nabla\log
p_t(\bm{x}|y) - \nabla\log p_t(\bm{x})$.
Sustituyendo en~\eqref{eq:classifier_guidance}:
\begin{align*}
  \nabla\log p_t^{(s)}(\bm{x}|y)
  &= \nabla\log p_t(\bm{x}) + s[\nabla\log p_t(\bm{x}|y) - \nabla\log p_t(\bm{x})]\\
  &= (1-s)\nabla\log p_t(\bm{x}) + s\,\nabla\log p_t(\bm{x}|y).
\end{align*}
Usando la identificación $\bm{s}_\theta = -\bm{\varepsilon}_\theta/\sigma_t$ y
multiplicando por $-\sigma_t$, se obtiene~\eqref{eq:cfg_eps}. $\square$
\end{proof}

La fórmula~\eqref{eq:cfg_eps} solo requiere dos evaluaciones de la misma red
(con y sin condición) por paso de muestreo: no es necesario ningún clasificador
externo. En la práctica, durante el entrenamiento se enmascara la condición $y$
con probabilidad $p_{\mathrm{drop}}\approx 0.1$, lo que entrena simultáneamente
el modelo condicionado y el no condicionado.

\begin{observacion}{El trade-off calidad-diversidad como temperatura}
El parámetro $s$ en CFG es exactamente la temperatura inversa de la distribución
de Gibbs del Capítulo~\ref{ch:cap5}. Para $s=1$: muestreo de $p_0(\bm{x}|y)$ exacto (en el
límite de red perfecta). Para $s>1$: concentración en regiones de alta
verosimilitud de la condición, imágenes más nítidas pero menos diversas. Para
$s<1$: imágenes más diversas pero menos fieles a la condición. Los valores
típicos son $s\in[1.5,\,10]$ según la tarea. El mismo trade-off entre calidad
de las muestras y diversidad que en el SGD bayesiano (Capítulo~\ref{ch:cap5}) reaparece aquí
como la elección del parámetro de guía.
\end{observacion}

\begin{fisicaconexion}{Temperatura, estadística bayesiana y el principio de máxima entropía}
La interpretación de temperatura del CFG conecta con un principio físico
fundamental. La distribución~\eqref{eq:guidance_distribución},
$p^{(s)}(\bm{x}|y) \propto p(\bm{x})\,p(y|\bm{x})^s$, es la solución del
problema de máxima entropía de Jaynes: maximizar la entropía de Shannon
$H[p] = -\int p\log p$ sujeto a la restricción de que la log-verosimilitud
media de la condición sea $\langle\log p(y|\bm{x})\rangle = C_s$ para alguna
constante $C_s$ que crece con $s$. El multiplicador de Lagrange de esa
restricción es exactamente $s$ (temperatura inversa), recuperando la forma de
Gibbs $p^{(s)} \propto p(\bm{x}) e^{s\log p(y|\bm{x})}$.

En \textbf{física estadística cuántica}, el mismo mecanismo aparece en el
ensemble de Gibbs con temperatura variable: reducir $T$ (aumentar $\beta = 1/k_BT$)
concentra la distribución en los estados de mínima energía, perdiendo diversidad
configuracional pero ganando fidelidad al estado de menor energía. El ``recocido
simulado'' del Capítulo~\ref{ch:cap5} es el proceso inverso: calentar para escapar
mínimos locales, enfriar para converger al mínimo global. El CFG con $s$ alto es
el recocido simulado del espacio de imágenes: ``enfría'' la distribución posterior
para producir imágenes más probables dado el texto.

En espectroscopía, el parámetro $s$ del CFG tiene un análogo en la
\textbf{deconvolución regularizada}: estimar la señal verdadera $\bm{x}$ a
partir de la señal observada $\bm{y}$ con ruido requiere un prior $p(\bm{x})$
(análogo al modelo no condicionado) y una verosimilitud $p(\bm{y}|\bm{x})$
(análogo al clasificador). La regularización de Tikhonov del Capítulo~\ref{ch:cap2}
es el caso $s = \lambda/\sigma^2$: el parámetro de regularización es la
temperatura inversa del prior gaussiano. El CFG unifica la inferencia bayesiana
estándar ($s=1$), la regularización de Tikhonov ($s$ elegido por validación) y
el muestreo de máxima verosimilitud ($s\to\infty$) en un único parámetro continuo.
\end{fisicaconexion}
\section{La arquitectura: U-Net con atención y DiT}
\label{sec:arquitectura_difusion}
\dificultad{2}{2}
%───────────────────────────────────────────────────────────────────────────────

\subsection{Requerimientos funcionales}
\label{subsec:arq_requerimientos}

La red $\bm{\varepsilon}_\theta(\bm{x}_t, y, t)$ debe satisfacer cuatro
requerimientos simultáneos:
\begin{enumerate}
  \item \textbf{Multi-escala}: la difusión actúa de modo cualitativamente
        distinto a distintas frecuencias espaciales. Las frecuencias bajas
        (estructura global) se destruyen lentamente; las altas (textura, detalles)
        se destruyen rápido.
  \item \textbf{Condición de tiempo}: $t$ determina el nivel de ruido
        $\sigma_t = \sqrt{1-\bar{\alpha}_t}$. La red debe saber en qué punto
        del forward process se encuentra para calibrar la escala de la
        predicción.
  \item \textbf{Condición a la señal}: el texto, la clase o la imagen de
        referencia $y$ deben influir en la predicción de forma no local
        (una descripción textual afecta a toda la imagen, no a una región).
  \item \textbf{Invariancias}: la misma descripción debe producir imágenes
        plausibles independientemente de la traslación o rotación de los objetos.
\end{enumerate}

\subsection{U-Net con atención}
\label{subsec:unet_atencion}

El U-Net (Ronneberger et al., 2015) con bloques de atención (Dhariwal \& Nichol,
2021) satisface estos requerimientos mediante tres mecanismos:

\textbf{Encoder-decoder con skip connections:} la arquitectura en forma de U
procesa la imagen a resoluciones progresivamente más bajas (encoder) y luego la
reconstruye (decoder). Las skip connections concatenan los mapas de características
del encoder con los del decoder a cada resolución, preservando los detalles de
alta frecuencia. Cada nivel del encoder captura una escala diferente del ruido.

\textbf{Embedding de tiempo:} el paso de tiempo $t$ se convierte en un
embedding mediante la codificación sinusoidal del Capítulo~\ref{ch:cap9}:
\[
  \bm{e}_t = \mathrm{MLP}\!\left(\left[\sin(10^{k/d}\pi t),\,\cos(10^{k/d}\pi t)
  \right]_{k=0}^{d/2-1}\right).
\]
El embedding $\bm{e}_t \in \mathbb{R}^{d_e}$ se inyecta en cada bloque residual
mediante suma o multiplicación elemento a elemento con los mapas de características:
$\bm{h} \leftarrow \bm{h} + \bm{W}_e\bm{e}_t$. Esto modula la escala de la
predicción en función del nivel de ruido.

\textbf{Cross-attention para la condición:} el texto $y$ se codifica por un
encoder de lenguaje (CLIP, T5) en una secuencia de embeddings $\bm{C}\in
\mathbb{R}^{L\times d_c}$. Los bloques de atención en las resoluciones bajas
del bottleneck del U-Net usan cross-attention con $\bm{C}$ como claves y
valores:
\[
  \mathrm{Attn}(\bm{Q},\bm{K},\bm{V}) = \mathrm{softmax}\!\left(\frac{\bm{Q}\bm{K}^\top}
  {\sqrt{d_k}}\right)\bm{V}, \quad \bm{K} = \bm{W}_K\bm{C},\; \bm{V} = \bm{W}_V\bm{C}.
\]
Este mecanismo, introducido en el Capítulo~\ref{ch:cap9}, permite que cada región espacial
de la imagen ``consulte'' qué parte del texto es más relevante para ella.

\subsection{Diffusion Transformer (DiT)}
\label{subsec:dit}

El DiT (Peebles \& Xie, 2023) reemplaza el U-Net por un Vision Transformer
(ViT) que opera en parches del espacio latente. Para una imagen latente de
tamaño $H\times W$ con $C$ canales, se divide en $N = HW/p^2$ parches de tamaño
$p\times p$, cada uno aplanado a un vector $\bm{z}_i \in \mathbb{R}^{p^2 C}$
y proyectado a la dimensión del transformer $d$: $\bm{z}_i \mapsto \bm{W}_p\bm{z}_i
+ \bm{e}_i^{\mathrm{pos}}$.

La condición de tiempo y clase se inyecta mediante \textbf{adaptive layer norm}
(adaLN): en lugar de layer norm estándar con parámetros fijos $(\gamma,\beta)$,
los parámetros son funciones de $(t,y)$:
\[
  \mathrm{adaLN}(\bm{h},t,y)
  = \gamma(t,y)\cdot\frac{\bm{h}-\mu(\bm{h})}{\sigma(\bm{h})} + \beta(t,y),
\]
donde $\gamma(t,y)$ y $\beta(t,y)$ son producidos por una red pequeña (MLP de
2 capas) aplicada al embedding conjunto $(t,y)$. Este mecanismo es computacionalmente
eficiente: el costo de adaLN es despreciable respecto al costo de la atención.

\textbf{Ventaja del DiT sobre el U-Net:} la atención global del transformer
permite que cualquier parche de la imagen interaccione directamente con cualquier
otro en un solo paso de atención, sin depender de la resolución a la que se
procesa. El U-Net necesita las capas de bottleneck para capturar dependencias
de largo alcance; el DiT las captura en todas las capas. Para imágenes de alta
resolución ($512\times 512$, $1024\times 1024$), el DiT escala mejor porque la
complejidad cuadrática de la atención se aplica a los parches del espacio latente
(dimensión reducida), no a los píxeles.

%───────────────────────────────────────────────────────────────────────────────
\section{Difusión en espacio latente}
\label{sec:latent_diffusion}
\dificultad{3}{2}
%───────────────────────────────────────────────────────────────────────────────

\subsection{La motivación dimensional}
\label{subsec:latent_motivacion}

La difusión directa en el espacio de imágenes $512\times 512\times 3 =
786{,}432$ dimensiones es computacionalmente prohibitiva: el reverse SDE opera
en $\mathbb{R}^{786432}$ con evaluaciones de red de $10^3$ pasos. La solución
es comprimir primero con un autoencoder a un espacio latente $\bm{z}\in
\mathbb{R}^{64\times 64\times 4}$ (factor de reducción $\approx 48\times$) y
hacer la difusión en el latente.

\subsection{El encoder, el decoder y la distribución latente}
\label{subsec:latent_vae}

El autoencoder consiste en un encoder $\mathcal{E}: \mathbb{R}^{H\times W\times C}
\to \mathbb{R}^{h\times w\times c}$ y un decoder $\mathcal{D}: \mathbb{R}^{h\times
w\times c} \to \mathbb{R}^{H\times W\times C}$ con $\mathcal{D}(\mathcal{E}
(\bm{x})) \approx \bm{x}$. El encoder VQ-VAE o KL-reg VAE de Rombach et al.\
(2022) añade regularización a la distribución de los latentes para que
$p_{\bm{z}} = \mathcal{E}_\#p_0$ sea próxima a $\mathcal{N}(\bm{0},\bm{I})$:
esto facilita que el prior del proceso de difusión (también $\mathcal{N}(\bm{0},
\bm{I})$) sea una aproximación razonable a la distribución latente real.

\subsection{El score en el espacio latente: análisis del jacobiano}
\label{subsec:latent_score}

Sea $\bm{z} = \mathcal{E}(\bm{x})$ y $p_{\bm{z}} = \mathcal{E}_\#p_0$ la
distribución pushforward. El score del proceso de difusión en el espacio latente
es $\nabla_{\bm{z}}\log p_{\bm{z},t}(\bm{z})$.

\begin{proposicion}{Score latente como pullback del score original}
\label{prop:score_latente}
Si $\mathcal{E}$ es invertible (localmente) con jacobiano $\bm{J}_{\mathcal{E}}
(\bm{x}) = \partial\mathcal{E}/\partial\bm{x} \in \mathbb{R}^{(hwc)\times(HWC)}$,
entonces el score de $p_{\bm{z},0}$ en $\bm{z} = \mathcal{E}(\bm{x})$ es:
\begin{equation}
  \nabla_{\bm{z}}\log p_{\bm{z},0}(\bm{z})
  = \bm{J}_{\mathcal{E}}(\bm{x})^{-\top}\nabla_{\bm{x}}\log p_0(\bm{x}),
  \label{eq:score_latente}
\end{equation}
donde $\bm{J}_{\mathcal{E}}^{-\top}$ es la transpuesta de la inversa del
jacobiano del encoder.
\end{proposicion}

\begin{proof}
Por la fórmula de cambio de variables para distribuciones:
$p_{\bm{z},0}(\bm{z}) = p_0(\mathcal{E}^{-1}(\bm{z}))\,|\det\bm{J}_\mathcal{E}^{-1}|$.
Tomando logaritmo:
$\log p_{\bm{z},0}(\bm{z}) = \log p_0(\bm{x}) - \log|\det\bm{J}_\mathcal{E}(\bm{x})|$.
Tomando gradiente respecto a $\bm{z}$ y usando la regla de la cadena con
$\partial\bm{x}/\partial\bm{z} = \bm{J}_\mathcal{E}^{-1}(\bm{x})$:
\begin{align*}
  \nabla_{\bm{z}}\log p_{\bm{z},0}(\bm{z})
  &= \bm{J}_\mathcal{E}^{-\top}\nabla_{\bm{x}}\log p_0(\bm{x})
  - \nabla_{\bm{z}}\log|\det\bm{J}_\mathcal{E}(\bm{x})|.
\end{align*}
El segundo término es la \textbf{corrección por cambio de volumen}: es cero si
$\mathcal{E}$ preserva el volumen ($|\det\bm{J}_\mathcal{E}| = 1$) y no nulo
en general. Para encoders con regularización KL, la distribución latente es
próxima a $\mathcal{N}(\bm{0},\bm{I})$ y el término de corrección es pequeño
cuando la regularización es fuerte. En la práctica (LDM, Stable Diffusion),
se entrena la red de score directamente en $\bm{z}$, absorbiendo implícitamente
la corrección. $\square$
\end{proof}

\begin{observacion}{Qué ignora la LDM y por qué funciona}
La LDM (Latent Diffusion Model) de Rombach et al.\ ignora el término
$\nabla_{\bm{z}}\log|\det\bm{J}_\mathcal{E}|$ de~\eqref{eq:score_latente}.
Esta aproximación es válida cuando (1) el encoder es aproximadamente isométrico
(preserva distancias), o (2) el proceso de difusión en el latente es
suficientemente largo para que el score en $t$ pequeño esté dominado por la
geometría de $p_{\bm{z},0}$ y no por la corrección de volumen. La regularización
KL del VAE es precisamente la herramienta que acerca $p_{\bm{z},0}$ a
$\mathcal{N}(\bm{0},\bm{I})$, reduciendo el impacto de la corrección en $t$
grande. El precio es que el decoder introduce un error de reconstrucción
$\|\bm{x} - \mathcal{D}(\mathcal{E}(\bm{x}))\|$ que no puede ser corregido
por el proceso de difusión.
\end{observacion}

\begin{fisicaconexion}{Reducción dimensional y espacios latentes en física experimental}
La difusión en espacio latente del LDM replica una estrategia ubicua en física
experimental: comprimir datos de alta dimensión a una representación latente
de baja dimensión antes de modelar su distribución. En \textbf{física de
partículas}, los detectores del LHC producen $O(10^8)$ canales por evento, pero
la física relevante vive en un espacio de variables de mucho menor dimensión
(momento transverso, multiplicidad de jets, energía faltante). Los autoencoders
de la Sección~\ref{sec:latent_diffusion} aplicados a eventos del LHC comprimen
$\bm{x}\in\mathbb{R}^{10^5}$ a $\bm{z}\in\mathbb{R}^{64}$, preservando la
información física discriminante (análogo al encoder KL-reg que preserva la
geometría de $p_0$) y descartando el ruido del detector (análogo al error de
reconstrucción $\|\bm{x}-\mathcal{D}(\mathcal{E}(\bm{x}))\|$).

La corrección del jacobiano~\eqref{eq:score_latente} tiene una interpretación
directa en este contexto. Si el encoder no preserva volúmenes
($|\det\bm{J}_\mathcal{E}| \neq 1$), las distancias en el espacio latente no
reflejan fielmente las distancias en el espacio original: una región de alta
densidad en $p_0$ puede aparecer como baja densidad en $p_{\bm{z},0}$ si
$|\det\bm{J}_\mathcal{E}|$ es grande allí. En la práctica de HEP, el equivalente
es la \textbf{función de aceptancia} del detector: la eficiencia de reconstrucción
$\epsilon(\bm{x})$ actúa como el jacobiano, y ignorarla produce sesgos
sistemáticos en las distribuciones modeladas. La regularización KL del VAE es
el análogo de la corrección por aceptancia: fuerza la distribución latente hacia
$\mathcal{N}(\bm{0},\bm{I})$, reduciendo la sensibilidad a las inhomogeneidades
del encoder/detector.

Para la simulación de calorímetros (el contexto del IAN-R1 y del LHC), el
espacio latente de la difusión puede interpretarse como el espacio de
\textbf{variables de forma de lluvia}: la distribución lateral, la profundidad
media y las fluctuaciones de energía son variables de baja dimensión que
capturan la física esencial de la lluvia electromagnética. Los modelos
CaloFlow, CaloDiT y CaloScore operan en este espacio latente físicamente
motivado, logrando una aceleración de $10^4$--$10^6\times$ respecto a GEANT4
sin sacrificar la fidelidad en las observables de interés.
\end{fisicaconexion}
\label{subsec:latent_score_difusion}

Para el proceso de difusión en el latente, la distribución marginal es:
\[
  p_{\bm{z},t}(\bm{z})
  = \int p_{\bm{z},0}(\bm{z}_0)\,q_t(\bm{z}|\bm{z}_0)\,d\bm{z}_0,
\]
con kernel de transición $q_t(\bm{z}|\bm{z}_0) = \mathcal{N}(\bm{z};
\sqrt{\bar{\alpha}_t}\bm{z}_0, (1-\bar{\alpha}_t)\bm{I})$.
El score es:
\[
  \nabla_{\bm{z}}\log p_{\bm{z},t}(\bm{z})
  = \frac{\mathbb{E}_{\bm{z}_0|{\bm{z},t}}\!\left[
    \nabla_{\bm{z}}\log q_t(\bm{z}|\bm{z}_0)\,p_{\bm{z},0}(\bm{z}_0)\right]}
    {p_{\bm{z},t}(\bm{z})}
  = -\frac{\bm{z} - \sqrt{\bar{\alpha}_t}\mathbb{E}[\bm{z}_0|\bm{z},t]}
    {1-\bar{\alpha}_t}.
\]
Esta expresión es idéntica a la del proceso OU en el espacio original, pero
ahora $\bm{z}_0$ es el latente del dato. El predictor de ruido $\bm{\varepsilon}_\theta
(\bm{z}_t, y, t)$ en el espacio latente es funcionalmente idéntico al predictor
en el espacio de píxeles, pero opera en dimensión $hwc \ll HWC$.

%───────────────────────────────────────────────────────────────────────────────
\section{Resumen y cierre del libro}
\label{sec:resumen12}
%───────────────────────────────────────────────────────────────────────────────

\subsection{El arco completo}
\label{subsec:arco_completo}

El libro partió de la pregunta más elemental: ¿qué es una red neuronal y por
qué funciona? Y llegó, por una cadena de pasos matemáticos rigurosos, a entender
exactamente por qué Stable Diffusion genera imágenes coherentes y qué aproxima
cuando lo hace. El camino:

\begin{enumerate}
  \item \textbf{Caps.\ 1--4} (MLP, backpropagation, optimización): la red
        neuronal es un aproximador universal entrenado por descenso de gradiente.
        El gradiente se propaga hacia atrás por la regla de la cadena.

  \item \textbf{Cap.\ 5} (Langevin/SGD): el SGD con ruido de mini-batch es una
        discretización de la SDE de Langevin $d\bm{\theta} = -\nabla\mathcal{L}
        \,dt + \sigma\,d\bm{W}$. Su distribución estacionaria es la distribución
        de Gibbs $p^*\propto e^{-2\mathcal{L}/\sigma^2}$.

  \item \textbf{Cap.\ 6} (Pontryagin): la backpropagation es el caso discreto
        del principio del máximo de Pontryagin. El estado adjunto son los
        gradientes de la pérdida respecto a las activaciones.

  \item \textbf{Cap.\ 7} (Probabilidad bayesiana): la distribución de Gibbs
        del Cap.\ 5 es la distribución posterior en un modelo bayesiano con
        prior plano y verosimilitud $e^{-\mathcal{L}/(\sigma^2/2)}$.

  \item \textbf{Caps.\ 8--9} (MLP en práctica, arquitecturas): dropout,
        normalización, transformers, redes equivariantes. Las herramientas que
        hacen funcionar las redes en alta dimensión.

  \item \textbf{Cap.\ 10} (Itô): el cálculo que hace rigurosas las SDEs. La
        variación cuadrática $[W,W]_T = T$ explica por qué $(dW)^2 = dt$ y por
        qué la fórmula de cambio de variables tiene un término de corrección de
        segundo orden.

  \item \textbf{Cap.\ 11} (Fokker-Planck, Anderson): la ecuación de Fokker-Planck
        conecta las SDEs con las EDPs. El teorema de Anderson determina la SDE
        del proceso reverso: el score $\nabla\log p_t$ es el único ingrediente
        desconocido. El denoising score matching convierte eso en una pérdida
        computable.

  \item \textbf{Cap.\ 12} (Modelos de difusión): DDPM, score SDE, DDIM, flow
        matching y CFG son instancias del mismo problema variacional con distintas
        elecciones de discretización y regularización. La arquitectura U-Net/DiT
        en el espacio latente implementa la función de score a escala.
\end{enumerate}

\subsection{La ecuación que sintetiza el libro}
\label{subsec:ecuacion_central}

Hay una ecuación que aparece en forma distinta en tres capítulos diferentes.

En el \textbf{Capítulo~\ref{ch:cap5}}, como SDE de Langevin del SGD:
\[
  d\bm{\theta}_t = -\nabla_{\bm{\theta}}\mathcal{L}(\bm{\theta}_t)\,dt
  + \sqrt{\eta c}\,d\bm{W}_t.
\]
El drift destruye correlaciones (escapa mínimos locales). El proceso converge
a la distribución de Gibbs $p^* \propto e^{-2\mathcal{L}/\eta c}$.

En el \textbf{Capítulo~\ref{ch:cap11}}, como reverse SDE de Anderson:
\[
  d\tilde{\bm{X}}_t
  = \left[\tilde{\bm{X}}_t + 2\nabla_{\bm{x}}\log p_{T-t}(\tilde{\bm{X}}_t)
  \right]dt + \sqrt{2}\,d\tilde{\bm{W}}_t.
\]
El drift reconstruye correlaciones (genera datos desde ruido). El proceso lleva
$\mathcal{N}(\bm{0},\bm{I})$ hacia $p_0$.

En el \textbf{Capítulo~\ref{ch:cap12}}, como pérdida de entrenamiento:
\[
  L_{\mathrm{DDPM}}(\theta)
  = \mathbb{E}\|\bm{\varepsilon} - \bm{\varepsilon}_\theta(\sqrt{\bar{\alpha}_t}
  \bm{x}_0 + \sqrt{1-\bar{\alpha}_t}\bm{\varepsilon}, t)\|^2.
\]
La red aprende el score que invierte la destrucción de correlaciones.

Las tres son la misma ecuación. La SDE de Langevin del SGD y el reverse SDE de
Anderson difieren solo en el signo del drift y en la interpretación del tiempo.
Entrenar una red de difusión es encontrar el campo vectorial que invierte la
dinámica del SGD. Generar una imagen es correr el SGD al revés.

\subsection{Conexiones abiertas}
\label{subsec:conexiones_abiertas}

El libro cierra el ciclo de los modelos de difusión, pero deja abiertas tres
conexiones hacia matemáticas más ricas que el lector interesado puede explorar:

\textbf{Ecuaciones de Hamilton-Jacobi y el control estocástico.} El score
$\nabla\log p_t$ es la solución de la ecuación de Hamilton-Jacobi estocástica
asociada al problema variacional del puente de Schrödinger. La conexión entre
modelos de difusión y control óptimo estocástico es un campo activo de
investigación que extiende el Capítulo~\ref{ch:cap6} al caso estocástico.

\textbf{Geometría de información y la métrica de Fisher.} El score $\nabla\log p$
es el gradiente en la métrica de Fisher-Rao sobre el espacio de distribuciones.
El descenso de gradiente natural usa la métrica de Fisher como tensor de
preacondicionamiento; el flujo de difusión en la métrica de Fisher produce
géodesicas en el espacio de distribuciones.

\textbf{Física de reactores y difusión heterogénea.} El término $\nabla\cdot\bm{D}$
del teorema de Anderson generalizado (Sección~11.3) tiene implicaciones directas
para la simulación de transporte de neutrones en medios heterogéneos. El proceso
reverso con difusividad variable permite muestrear distribuciones de flujo
neutrónico con la misma maquinaria de los modelos generativos.

%───────────────────────────────────────────────────────────────────────────────
\label{sec:estrella12}
\begin{estrella}{Modelos generativos para simulación de calorímetros}

\textbf{El problema: GEANT4 como cuello de botella computacional.}
La simulación de lluvia de partículas en el calorímetro electromagnético del
LHC (ATLAS, CMS) con GEANT4 requiere $O(10^3)$ segundos de CPU por evento a
alta energía. Para los $10^9$ eventos de la Run 3, la simulación completa
consumiría $O(10^{12})$ segundos: inviable sin aproximaciones. Los modelos
generativos del Capítulo~\ref{ch:cap12} son la alternativa que el CERN explora
sistemáticamente para reemplazar GEANT4 con emuladores $10^4$--$10^6\times$ más
rápidos.

\textbf{Arquitectura de los modelos de calorímetro.}
El problema tiene una estructura natural que motiva el diseño del modelo:
\begin{itemize}
  \item \textbf{Forward process}: el ruido del forward OU actúa sobre el mapa
        de energía $\bm{x}\in\mathbb{R}^{N_\eta\times N_\phi\times N_r}$ (cuadrícula
        de vóxeles $\eta$-$\phi$-radio), llevándolo hacia $\mathcal{N}(\bm{0},\bm{I})$.
  \item \textbf{Arquitectura}: la simetría de traslación en $\phi$ (el calorímetro
        es azimutalmente simétrico) justifica el uso de capas convolucionales
        equivariantes (Sección~\ref{sec:cnn}) en $\phi$; la no uniformidad en
        $\eta$ y $r$ favorece la atención global (DiT, Sección~\ref{subsec:dit}).
  \item \textbf{Condicionamiento}: la energía incidente $E_{\mathrm{inc}}$ y el
        tipo de partícula $\pi/e/\gamma$ se inyectan como condición $y$ vía CFG
        o cross-attention, produciendo lluvias condicionadas físicamente correctas.
\end{itemize}

\textbf{Modelos específicos y su posición en el marco unificador.}
La Tabla~\ref{tab:unificacion} organiza los modelos de calorímetro según $\epsilon$:
CaloScore (Song et al., adaptado) usa el score SDE; CaloDiT y CaloFlow usan el
DiT con flow matching (OT, $\epsilon\to 0$); FastCaloGAN usa GANs (fuera del
marco de Schrödinger). La transición de DDPM a flow matching reduce los pasos
de muestreo de $10^3$ a $O(10)$, lo que es crítico para la producción: a
$10^9$ eventos, cada paso de evaluación de red adicional cuesta horas de CPU.

\textbf{El reverse SDE con difusividad variable para el IAN-R1.}
Para el reactor de investigación IAN-R1, el flujo neutrónico $\Phi(\bm{r})$
tiene simetría cilíndrica $\mathrm{SO}(2)\times\mathbb{Z}_2$ y difusividad
$D(\bm{r})$ variable entre el combustible y el reflector de berilio (factor
$\sim 6$). El reverse SDE general~\eqref{eq:reverse_general} del
Teorema~\ref{thm:anderson_general}, con el término $\nabla\cdot\bm{D}$ no nulo,
genera distribuciones de flujo que respetan la heterogeneidad del medio:
\begin{itemize}
  \item El forward process lleva $\Phi(\bm{r})$ hacia $\mathcal{N}(\bm{0},\bm{I})$
        en tiempo $T$ suficientemente largo (comparable al tiempo de difusión
        de neutrones sobre la escala del núcleo).
  \item El score network $\bm{s}_\theta(\bm{z}_t, E_{\mathrm{barras}}, t)$ se
        entrena con $N\sim 10^4$ simulaciones de MCNP o difusión de dos grupos,
        condicionado a la configuración de barras de control
        $E_{\mathrm{barras}}\in\mathbb{R}^d$.
  \item El muestreo con flow matching (trayectorias rectas, $O(10)$ pasos)
        genera en milisegundos distribuciones de flujo para nuevas configuraciones
        de barras, reemplazando horas de MCNP.
\end{itemize}
El modelo completo es un emulador físicamente consistente del reactor:
respeta la simetría cilíndrica (via CNN equivariante en $\phi$), la
heterogeneidad del medio (via $\nabla\cdot\bm{D}\neq\bm{0}$ en el reverse SDE),
y las condiciones de contorno (flujo cero en la frontera del reflector,
implementada como restricción en el espacio latente del encoder).
\end{estrella}

%───────────────────────────────────────────────────────────────────────────────
\section*{Ejercicios computacionales}
\addcontentsline{toc}{section}{Ejercicios computacionales}
%───────────────────────────────────────────────────────────────────────────────

\begin{ejerciciocomp}{Implementación completa de DDPM y DDIM en 1D}
\textbf{Objetivo:} implementar el Algoritmo~\ref{alg:ddpm} completo y verificar
que DDIM con $\eta=0$ genera muestras de igual calidad con menos pasos.

\medskip\noindent\textbf{Algoritmo:}
\begin{enumerate}
  \item Usar $p_0 = \frac{1}{2}\mathcal{N}(-2,0.3^2)+\frac{1}{2}\mathcal{N}(2,0.3^2)$
        (bimodal). Implementar el forward process con $T=500$ pasos y schedule
        coseno. Verificar que $q(\bm{x}_T|\bm{x}_0)\approx\mathcal{N}(0,1)$.
  \item Entrenar un MLP $(2, 128, 128, 128, 1)$ con entradas $(x_t, t_{\mathrm{embed}})$
        y pérdida~\eqref{eq:ddpm_loss}, con $t$ inyectado como embedding
        sinusoidal de dimensión 32. Usar Adam con $\eta=3\times10^{-4}$,
        $5\times10^4$ iteraciones.
  \item Muestrear $10^3$ puntos con DDPM ($T=500$ pasos) y con DDIM ($\eta=0$)
        para $N\in\{5,10,20,50,100,500\}$ pasos. Calcular la distancia de
        Wasserstein $W_1$ entre las muestras y $p_0$ en cada caso.
  \item Verificar el Teorema~\ref{thm:ddim}: que DDIM con $\eta=0$ produce
        trayectorias deterministas (misma muestra inicial $\to$ misma imagen
        final) y DDIM con $\eta=1$ produce trayectorias estocásticas.
\end{enumerate}
\medskip\noindent\textbf{Verificación:} DDIM con $N=50$ debe alcanzar
$W_1 < 0.05$; DDPM necesita $N\geq 200$ para el mismo umbral.
\end{ejerciciocomp}

\begin{ejerciciocomp}{Classifier-Free Guidance: verificación de la interpretación de temperatura}
\textbf{Objetivo:} verificar empíricamente la Proposición~\ref{prop:guidance_temperatura}
— que CFG con escala $s$ muestrea de $p(\bm{x})\,p(y|\bm{x})^s$.

\medskip\noindent\textbf{Algoritmo:}
\begin{enumerate}
  \item Usar $p_0(\bm{x}|y=1)=\mathcal{N}(-3,0.5^2)$ y $p_0(\bm{x}|y=2)=
        \mathcal{N}(3,0.5^2)$ con prior uniforme $p(y=1)=p(y=2)=0.5$.
        Entrenar un modelo DDPM condicionado ($T=200$, mismo MLP del ejercicio
        anterior), enmascarando $y$ con probabilidad $0.1$.
  \item Para $s\in\{0.5, 1, 2, 5, 10\}$, generar $2\times10^3$ muestras con
        CFG para $y=1$ usando la fórmula~\eqref{eq:cfg_eps}.
  \item Para cada $s$, estimar la distribución de muestras con KDE. Verificar
        que la media y varianza se aproximan a $\mu_s\approx -3$ y
        $\sigma_s^2\approx 0.25/s$ (predicción de la
        Proposición~\ref{prop:guidance_temperatura}).
  \item Trazar la curva \textit{calidad vs.\ diversidad}: en el eje $x$,
        la fracción de muestras en $(-5,-0.5)$ (probabilidad de moda correcta);
        en el eje $y$, la varianza de las muestras. ¿En qué $s$ se maximiza
        un score $F_1$ balanceado entre ambos?
\end{enumerate}
\end{ejerciciocomp}

%───────────────────────────────────────────────────────────────────────────────
\begin{notasbib}
El modelo DDPM fue introducido por Ho et al.\ (2020)~\cite{ho2020denoising},
citado en el Capítulo~\ref{ch:cap11}. El noise schedule coseno es de Nichol y
Dhariwal (2021)~\cite{nichol2021improved}. El marco unificado score-based SDE,
incluyendo la VP-SDE y VE-SDE, es de Song et al.\ (2021)~\cite{song2021score},
citado en el Capítulo~\ref{ch:cap11}. El sampler DPM-Solver es de Lu et al.\
(2022)~\cite{lu2022dpm}.

DDIM fue introducido por Song, Meng y Ermon (2020)~\cite{song2020denoising}. La
demostración del Teorema~\ref{thm:ddim} presentada aquí, mostrando que DDIM con
$\eta=0$ es la discretización de Euler de la ODE de flujo de probabilidad del
Capítulo~\ref{ch:cap11}, sigue la presentación de Karras et al.\
(2022)~\cite{karras2022elucidating}, que también unifica múltiples samplers bajo
un mismo marco.

El flow matching y el rectified flow son de Lipman et al.\
(2022)~\cite{lipman2022flow} y Liu et al.\ (2022)~\cite{liu2022flow},
ambos citados en el Capítulo~\ref{ch:cap11}. El OT flow matching es de Tong et
al.\ (2023)~\cite{tong2023improving}, también citado allí.

El classifier guidance fue introducido por Dhariwal y Nichol
(2021)~\cite{dhariwal2021diffusion}. El Classifier-Free Guidance (CFG) fue
introducido por Ho y Salimans (2022)~\cite{ho2022classifier}. La interpretación
de CFG como temperatura del posterior bayesiano~\eqref{eq:guidance_distribución}
es estándar; la conexión con el principio de máxima entropía de Jaynes y la
regularización de Tikhonov presentada en la \textsf{Conexión física} sigue a
la perspectiva de Chung et al.\ (2022)~\cite{chung2022diffusion}.

El U-Net original es de Ronneberger, Fischer y Brox (2015)~\cite{ronneberger2015unet}.
La versión con bloques de atención para difusión es de Dhariwal y Nichol
(2021)~\cite{dhariwal2021diffusion}. El Diffusion Transformer (DiT) es de
Peebles y Xie (2023)~\cite{peebles2023scalable}; la adaptive layer norm fue
introducida en ese trabajo. El mecanismo de cross-attention para condicionamiento
por texto es de Rombach et al.\ (2022)~\cite{rombach2022high}, que también
introdujo el Latent Diffusion Model (LDM) y Stable Diffusion. La Proposición
sobre el score latente como pullback~\eqref{eq:score_latente} sigue a la
discusión de Pidstrigach (2022)~\cite{pidstrigach2022score}.

Para la aplicación a simulación de calorímetros (Sección~\ref{sec:estrella12}),
los modelos de referencia son CaloFlow (Kruse et al., 2021), CaloScore (Mikuni
y Nachman, 2022)~\cite{mikuni2022score}, y CaloDiT (Acosta et al., 2023). Una
revisión comparativa de los modelos generativos para calorímetros es Paganini et
al.\ (2018)~\cite{paganini2018calogan} y el Fast Calorimeter Simulation Challenge
(Faucci Giannelli et al., 2022)~\cite{faucci2022fast}.
\end{notasbib}

%───────────────────────────────────────────────────────────────────────────────
\section*{Ejercicios resueltos}
\addcontentsline{toc}{section}{Ejercicios resueltos}
%───────────────────────────────────────────────────────────────────────────────

\begin{ejercicio}
\label{ej:ddpm_posterior}
\textbf{(Verificación del posterior condicionado a $\bm{x}_0$).}
Verificar que $q(\bm{x}_{t-1}|\bm{x}_t,\bm{x}_0)$ es gaussiana con la media
y varianza de la Proposición~\ref{prop:ddpm_posterior}, y que cuando $\beta_t\to 0$
la varianza $\tilde{\beta}_t\to 0$ y la media $\tilde{\bm{\mu}}_t\to\bm{x}_{t-1}$
(el proceso es casi determinista para pasos pequeños).

\medskip\noindent\textit{Solución.}
Aplicando Bayes: $q(\bm{x}_{t-1}|\bm{x}_t,\bm{x}_0)\propto q(\bm{x}_t|\bm{x}_{t-1})
q(\bm{x}_{t-1}|\bm{x}_0)$. Las densidades son:
\[
  -2\log q(\bm{x}_t|\bm{x}_{t-1})
  = \frac{\|\bm{x}_t-\sqrt{\alpha_t}\bm{x}_{t-1}\|^2}{\beta_t} + \text{cte},
  \quad
  -2\log q(\bm{x}_{t-1}|\bm{x}_0)
  = \frac{\|\bm{x}_{t-1}-\sqrt{\bar{\alpha}_{t-1}}\bm{x}_0\|^2}
    {1-\bar{\alpha}_{t-1}} + \text{cte}.
\]
La suma es cuadrática en $\bm{x}_{t-1}$ con coeficiente de $\|\bm{x}_{t-1}\|^2$:
$\frac{\alpha_t}{\beta_t} + \frac{1}{1-\bar{\alpha}_{t-1}} = \frac{\alpha_t(1-\bar{\alpha}_{t-1})
+\beta_t}{\beta_t(1-\bar{\alpha}_{t-1})} = \frac{1-\bar{\alpha}_t}{\beta_t(1-\bar{\alpha}_{t-1})}
= \frac{1}{\tilde{\beta}_t}$.
El coeficiente de $\bm{x}_{t-1}$ da la media $\tilde{\bm{\mu}}_t$.
Para $\beta_t\to 0$: $\alpha_t\to 1$, $\bar{\alpha}_{t-1}\approx\bar{\alpha}_t$,
luego $\tilde{\beta}_t = \beta_t(1-\bar{\alpha}_{t-1})/(1-\bar{\alpha}_t)\approx\beta_t\to 0$
y $\tilde{\bm{\mu}}_t \approx \bm{x}_{t-1}$. $\square$
\end{ejercicio}

\begin{ejercicio}
\label{ej:ddim_determinismo}
\textbf{(DDIM con $\eta=0$ como ODE de flujo).}
Para el VP-SDE con $\bar{\alpha}_t = e^{-t^2/2}$ (schedule cuadrático):
\textbf{(a)} calcular el paso DDIM con $\eta=0$ de $t=1.0$ a $t'=0.9$
para un $\bm{x}_t$ dado con predictor de ruido constante $\bm{\varepsilon}_\theta =
\bm{e}$ (vector unitario).
\textbf{(b)} Verificar que coincide con el paso Euler de la ODE de flujo.

\medskip\noindent\textit{Solución.}
\textbf{(a)} $\bar{\alpha}_1 = e^{-1/2}\approx 0.6065$, $\bar{\alpha}_{0.9} =
e^{-0.405}\approx 0.6670$. Paso DDIM:
\[
  \bm{x}_{0.9} = \sqrt{0.667}\hat{\bm{x}}_0 + \sqrt{1-0.667}\,\bm{e},
  \quad \hat{\bm{x}}_0 = \frac{\bm{x}_1 - \sqrt{1-0.6065}\,\bm{e}}{\sqrt{0.6065}}.
\]

\textbf{(b)} ODE de flujo VP: $d\bm{x}/dt = -\frac{\beta(t)}{2}\bm{x}_t
+ \frac{\beta(t)}{2\sqrt{1-\bar{\alpha}_t}}\bm{\varepsilon}_\theta$
con $\beta(t) = t$ (derivada del schedule cuadrático). Euler de $t=1$ a $0.9$:
$\bm{x}_{0.9} = \bm{x}_1 + 0.1\cdot(-\frac{1}{2}\bm{x}_1 + \frac{1}{2\sqrt{0.3935}}
\bm{e})$. Verificando numéricamente: ambos valores coinciden hasta $O((\Delta t)^2)$.
$\square$
\end{ejercicio}

%───────────────────────────────────────────────────────────────────────────────
\section*{Ejercicios propuestos}
\addcontentsline{toc}{section}{Ejercicios propuestos}
%───────────────────────────────────────────────────────────────────────────────

\begin{ejercicio}
\textbf{(Implementación completa de DDPM en 1D).}
Sea $p_0 = \frac{1}{3}\mathcal{N}(-3,0.5) + \frac{1}{3}\mathcal{N}(0,0.5)
+ \frac{1}{3}\mathcal{N}(3,0.5)$ (mezcla trimodal).
\begin{enumerate}[label=(\alph*)]
  \item Implementar el forward process con $T=200$ pasos y schedule lineal.
        Verificar que $q(\bm{x}_T|\bm{x}_0)\approx\mathcal{N}(0,1)$ para
        $T=200$.
  \item Entrenar un MLP de 3 capas con $256$ neuronas para predecir
        $\bm{\varepsilon}$ con la pérdida~\eqref{eq:ddpm_loss}.
        El tiempo $t$ se inyecta como embedding sinusoidal.
  \item Implementar el muestreo del Algoritmo~\ref{alg:ddpm}. Comparar
        el histograma de $10^4$ muestras con $p_0$.
  \item Repetir con schedule coseno. ¿Mejora la calidad de las muestras?
\end{enumerate}
\end{ejercicio}

\begin{ejercicio}
\textbf{(DDIM vs.\ DDPM: calidad en función del número de pasos).}
Usando la red $\bm{\varepsilon}_\theta$ entrenada en el ejercicio anterior:
\begin{enumerate}[label=(\alph*)]
  \item Implementar el paso DDIM~\eqref{eq:ddim_step} con $\eta=0$ y $\eta=1$.
  \item Para $N\in\{5, 10, 20, 50, 200\}$ pasos, comparar la distancia de
        Wasserstein $W_1$ entre las muestras y $p_0$ para DDPM, DDIM ($\eta=0$)
        y DDIM ($\eta=1$).
  \item ¿Para qué valor de $N$ DDIM ($\eta=0$) alcanza la misma calidad que
        DDPM con $N=200$?
  \item Graficar las trayectorias de muestreo para $5$ puntos: ¿son
        deterministas con $\eta=0$?
\end{enumerate}
\end{ejercicio}

\begin{ejercicio}
\textbf{(Classifier-Free Guidance en una mezcla bimodal condicionada).}
Sea $p_0(x|y=1) = \mathcal{N}(-3,0.5)$ y $p_0(x|y=2) = \mathcal{N}(3,0.5)$
con $p_0(y=1) = p_0(y=2) = 0.5$.
\begin{enumerate}[label=(\alph*)]
  \item Entrenar un modelo DDPM condicionado a $y$, enmascarando $y$ con
        probabilidad $0.1$ durante el entrenamiento.
  \item Para $s\in\{0.5, 1, 2, 5, 10\}$, generar $10^3$ muestras con CFG
        para $y=1$. Calcular $\mathbb{P}(\text{muestra}\in[-5,0])$ para cada $s$.
  \item Verificar la interpretación de temperatura: mostrar que la distribución
        de muestras se aproxima a $\mathcal{N}(-3, 0.5/s)$ para $s$ grande.
  \item ¿Para qué valor de $s$ el trade-off calidad-diversidad es óptimo
        según la métrica FID (Fréchet Inception Distance) discreta?
\end{enumerate}
\end{ejercicio}

\begin{ejercicio}
\textbf{(Difusión con difusividad variable: el caso del reactor esférico).}
Retomar el ejercicio del reactor esférico del Capítulo~\ref{ch:cap11} (Ejercicio propuesto 4)
con $D(r) = D_0(1-r^2/R^2)$ y $p^*(r)$ la distribución estacionaria calculada.
\begin{enumerate}[label=(\alph*)]
  \item Usar el reverse SDE del Teorema~\ref{thm:anderson_general} para generar
        muestras de $p^*(r)$ partiendo de $p_T = \mathcal{N}(0,\sigma^2)$.
        Verificar que la distribución generada coincide con $p^*(r)$.
  \item Entrenar una red para predecir el score $\nabla\log p_t^*(r)$ usando
        denoising score matching con el kernel de transición del forward process
        con $D(r)$.
  \item Comparar el reverse SDE con $D(r)$ variable (que incluye $\nabla\cdot\bm{D}
        = 2D_0(-2r/R^2) = -4D_0 r/R^2$) con el reverse SDE isótropo (que ignora
        ese término). ¿Qué error introduce ignorar $\nabla\cdot\bm{D}$?
  \item Interpretar el resultado en términos del transporte de neutrones: el
        forward process simula la dispersión de neutrones en el reactor, y el
        reverse process los ``enfoca'' de vuelta a la distribución de equilibrio.
\end{enumerate}
\end{ejercicio}
